\section{Seleção de Variáveis}

\subsection{Seleção Interna de Variáveis}

\begin{frame}

A seleção das variáveis é parte intríseca do treino do modelo.

\end{frame}

\begin{frame}

\frametitle{Árvore de Decisão}

\begin{itemize}
    \item Para cada variável, avalia o ganho de informação em diferentes partições
    \item Partições com maior ganho ficam mais próximas da raiz
\end{itemize}

\vspace{1cm}

\textbf{Principais problemas}
\begin{itemize}
    \item Pequenas mudanças nos dados podem levar a estruturas completamente diferentes - instabilidade
    \item Erros são fortemente cumulativos
\end{itemize}

\end{frame}

\begin{frame}

\frametitle{Floresta Aleatória}

Agrega multiplas árvores de decisão, com algumas modificações
\begin{itemize}
    \item Cada árvore é construída sobre um subconjunto dos dados
    \item Na construção de cada nó, um subconjunto das variáveis é considerado
    \item Cada árvore vota conforme seu resultado e a classificação é feita a partir destes
\end{itemize}

\vspace{1cm}

\textbf{Principais problemas}
\begin{itemize}
    \item Mitiga, mas não resolve a instabilidade
    \item Aumenta o custo computacional
\end{itemize}

\end{frame}

\subsection{Filtragem de Variáveis}

\begin{frame}

Analisa, estatisticamente, a relação entre cada variável e as classes do problema. Mantém aquelas que são fortes indicadores de classe.

\textbf{Exemplo}: Retirar variáveis cuja variança seja baixa no conjunto completo de dados, visto que estas tendem a não ser relevantes para a classificação.

\textbf{Principais problemas}
Por não considerar a relação entre variáveis,
\begin{itemize}
    \item pode remover aquelas que participam de relações úteis para a classificação
    \item pode manter variáveis redundantes
\end{itemize}
\end{frame}

\subsection{Métodos Encapsuladores}
\begin{frame}
    Encapsula o modelo em um problema de optimização.
    \begin{itemize}
        \item Seleciona um subconjunto das variáveis
        \item Treina o modelo sobre este subconjunto
        \item Avalia o modelo conforme uma métrica apropriada
        \item Seleciona o modelo com melhor performance
    \end{itemize}

\textbf{Principal problema}: custo computacional elevado
\end{frame}


